% ============================================================================
% Appendix: Effective Potential Along a Chosen Collective Path
% ============================================================================
% Status: [Dc] geometric profile within chosen-path displacement model
% Scope: This is NOT a derivation of V(q) from the full 5D action.
%        It is a geometric construction along a chosen deformation path.
% Phase 1 / R3 Mode B deliverable (see PHASE1_PLAN_REVISED.md v3.1)
% ============================================================================

\chapter{Effective Potential Along a Chosen Collective Path}
\label{app:vq_chosen_path}

\section*{Scope and Epistemic Status}

This appendix constructs the \textbf{geometric brane-energy contribution}
$V_{\mathrm{geom}}(q)$ to the effective potential along a chosen collective
deformation path for the Y-junction. It does \emph{not}:
\begin{itemize}
    \item derive the full effective potential $V(q)$ from the 5D action,
    \item complete the Put~C reduction corridor (steps C2--C4 remain open),
    \item establish the existence or stability of a secondary minimum,
    \item close the identification $\Earm \equiv \Delta m_{np}$.
\end{itemize}

The collective coordinate $q$ is a \textbf{chosen interpolation parameter},
not a uniquely derived normal mode of the full 5D dynamics. Different path
choices yield quantitatively different profiles. The qualitative feature---a
single-well geometric minimum at the Steiner configuration---is
path-independent.

\textbf{Intended epistemic tag:} [Dc] for the geometric profile within the
chosen displacement model. Full $V(q)$ closure requires completing Put~C
(5D action $\to$ effective 1D action), which is beyond the scope of this
appendix.

% ============================================================================
\section{Path Definition}
\label{sec:vq:path}
% ============================================================================

\subsection{Junction Node Displacement Model}

Following the Put~C corridor definition
(Chapter~\ref{ch:instanton}), the collective coordinate $q$ parameterizes
the displacement of the Y-junction node from its equilibrium (Steiner)
position. Two natural displacement directions exist:

\begin{enumerate}
    \item \textbf{In-brane displacement:} The node moves within the brane
          plane, away from the Steiner center toward one boundary point. This
          breaks $\Zthree$ to $\Ztwo$.
    \item \textbf{Compact-direction displacement:} The node moves along the
          compact coordinate $\xi$, as defined in the Put~C ansatz
          ($q = \xi_{\mathrm{node}} - \xi_0$). This preserves the in-brane
          $\Zthree$ symmetry.
\end{enumerate}

Both are \textbf{chosen paths}---neither is uniquely derived as the dominant
deformation mode. This appendix computes $V_{\mathrm{geom}}(q)$ for both
paths. The in-brane path is the primary model (it produces the
$\Zthree$-breaking geometry relevant to the metastable junction); the
compact-direction path is presented for comparison.

\subsection{Boundary Conditions}

The Y-junction has three worldsheet arms terminating at boundary points
$\{\mathbf{P}_1, \mathbf{P}_2, \mathbf{P}_3\}$ on the brane. These form
an equilateral triangle at distance $R$ from the Steiner center:
\begin{equation}
    \mathbf{P}_1 = (R, 0), \quad
    \mathbf{P}_2 = R\bigl(-\tfrac{1}{2}, \tfrac{\sqrt{3}}{2}\bigr), \quad
    \mathbf{P}_3 = R\bigl(-\tfrac{1}{2}, -\tfrac{\sqrt{3}}{2}\bigr)
    \label{eq:boundary_points}
    \tagDef
\end{equation}
where $R$ is the characteristic arm length. The Steiner point (center)
is at the origin. \tagDer{} (Ch.~\ref{ch:anchor}, Steiner optimality).

\subsection{Non-Uniqueness of Path Choice}

The displacement direction is chosen, not derived. For the in-brane path,
the $\Zthree$ symmetry of the boundary configuration means that choosing
the displacement toward any of the three vertices $\mathbf{P}_i$ gives an
equivalent profile (related by $\Zthree$ rotation). The path is unique up
to this $\Zthree$ equivalence.

Other deformation paths (e.g., angular distortion of arm directions,
combined in-brane and compact displacement) would yield different
$V_{\mathrm{geom}}(q)$ profiles. The Steiner minimum at $q = 0$ is
preserved in all cases (it is a property of the boundary conditions,
not the path choice), but the curvature, barrier shape, and large-$q$
behavior are path-dependent.

% ============================================================================
\section{Geometric Brane Energy: In-Brane Path}
\label{sec:vq:inplane}
% ============================================================================

\subsection{Arm-Length Formula}

The junction node is displaced by $q$ from the Steiner center toward
$\mathbf{P}_1$, placing the node at $(q, 0)$ with $0 \leq q < R$. The
arm lengths are:

\begin{derivationbox}[title=Arm Lengths Along In-Brane Path]
    \begin{align}
        L_1(q) &= R - q
            \label{eq:L1_inplane} \\
        L_2(q) = L_3(q) &= \sqrt{R^2 + Rq + q^2}
            \label{eq:L23_inplane}
    \end{align}
    \tagDc
\end{derivationbox}

\begin{proof}
$L_1$ is the distance from $(q, 0)$ to $(R, 0)$, giving $R - q$.
$L_2$ is the distance from $(q, 0)$ to $(-R/2, R\sqrt{3}/2)$:
\begin{equation*}
    L_2 = \sqrt{\bigl(q + \tfrac{R}{2}\bigr)^2 + \tfrac{3R^2}{4}}
        = \sqrt{q^2 + Rq + \tfrac{R^2}{4} + \tfrac{3R^2}{4}}
        = \sqrt{R^2 + Rq + q^2}
\end{equation*}
$L_3 = L_2$ by the residual $\Ztwo$ symmetry ($\mathbf{P}_2 \leftrightarrow \mathbf{P}_3$).
\end{proof}

\subsection{Total Arm Length and Geometric Energy}

The total arm length along the in-brane path:
\begin{equation}
    L_{\mathrm{tot}}(q) = (R - q) + 2\sqrt{R^2 + Rq + q^2}
    \label{eq:Ltot_inplane}
\end{equation}

The geometric brane energy (Nambu--Goto, Lemma~2 of Ch.~\ref{ch:anchor}):
\begin{equation}
    \boxed{V_{\mathrm{geom}}(q) = \tau \times L_{\mathrm{tot}}(q)}
    \label{eq:Vgeom_inplane}
    \tagDc
\end{equation}
where $\tau$ is the string/edge tension. The energy above the Steiner minimum:
\begin{equation}
    \Delta V_{\mathrm{geom}}(q) = \tau \bigl[L_{\mathrm{tot}}(q) - 3R\bigr]
    \label{eq:dVgeom}
\end{equation}

\subsection{Steiner Minimum}

\begin{theorem}[Steiner Minimum at $q = 0$]
\label{thm:steiner_min}
$V_{\mathrm{geom}}(q)$ has a strict local minimum at $q = 0$ with:
\begin{align}
    V_{\mathrm{geom}}'(0) &= 0 \label{eq:Vgeom_first} \\
    V_{\mathrm{geom}}''(0) &= \frac{3\tau}{2R} > 0 \label{eq:Vgeom_second}
\end{align}
\end{theorem}

\begin{proof}
Differentiating Eq.~\eqref{eq:Ltot_inplane}:
\begin{equation}
    \frac{dL_{\mathrm{tot}}}{dq} = -1 + \frac{R + 2q}{\sqrt{R^2 + Rq + q^2}}
\end{equation}
At $q = 0$: $dL/dq = -1 + R/R = 0$. \checkmark

For the second derivative:
\begin{equation}
    \frac{d^2 L_{\mathrm{tot}}}{dq^2}
    = \frac{3R^2/2}{\bigl(R^2 + Rq + q^2\bigr)^{3/2}}
\end{equation}
At $q = 0$: $d^2L/dq^2 = 3/(2R)$. Therefore
$V_{\mathrm{geom}}''(0) = \tau \times 3/(2R) = 3\tau/(2R) > 0$.
\end{proof}

This is a restatement of the Steiner optimality theorem
(Ch.~\ref{ch:anchor}) within the specific displacement model.

\subsection{Profile Shape}

The geometric energy rises monotonically from the Steiner minimum:
$\Delta V_{\mathrm{geom}}(q) > 0$ for all $q \neq 0$ within the domain
$0 < q < R$. There is no secondary minimum in $V_{\mathrm{geom}}$
alone---the Steiner point is the unique minimum of the brane
energy.

In dimensionless units ($\tilde{q} = q/R$):
\begin{equation}
    \frac{L_{\mathrm{tot}}(\tilde{q})}{R}
    = (1 - \tilde{q}) + 2\sqrt{1 + \tilde{q} + \tilde{q}^2}
    \label{eq:Ltot_dimless}
\end{equation}

\noindent Representative values:
\begin{center}
\begin{tabular}{lll}
\toprule
$\tilde{q}$ & $L_{\mathrm{tot}}/R$ & $\Delta V_{\mathrm{geom}} / (\tau R)$ \\
\midrule
$0.0$ & $3.000$ & $0.000$ \\
$0.1$ & $3.007$ & $0.007$ \\
$0.2$ & $3.027$ & $0.027$ \\
$0.5$ & $3.146$ & $0.146$ \\
$0.8$ & $3.477$ & $0.477$ \\
\bottomrule
\end{tabular}
\end{center}

% ============================================================================
\section{Geometric Brane Energy: Compact-Direction Path}
\label{sec:vq:compact}
% ============================================================================

For comparison, we compute $V_{\mathrm{geom}}$ for displacement along the
compact coordinate $\xi$ (the Put~C definition: $q = \xi_{\mathrm{node}} -
\xi_0$). The node moves to $(0, 0, \xi_0 + q)$ while all boundary points
remain at $(x_i, y_i, \xi_0)$ on the brane. By symmetry, all three arms
have equal length:
\begin{equation}
    L_i(q) = \sqrt{R^2 + q^2}, \quad
    L_{\mathrm{tot}}(q) = 3\sqrt{R^2 + q^2}
    \label{eq:Ltot_compact}
\end{equation}

Properties at $q = 0$:
\begin{align}
    L_{\mathrm{tot}}'(0) &= 0 \\
    V_{\mathrm{geom}}''(0) &= \frac{3\tau}{R}
\end{align}

The compact-direction path gives a stiffer potential (twice the curvature of
the in-brane path) and preserves $\Zthree$ for all $q$. Since the metastable
junction is interpreted as $\Zthree$-breaking
(Chapter~\ref{ch:metastable}~\S\ref{sec:metastable:z3}), the in-brane
path is the more physically relevant model for the anchor-to-metastable
deformation.

% ============================================================================
\section{Barrier Existence Along Chosen Path}
\label{sec:vq:barrier}
% ============================================================================

\subsection{Geometric Barrier Contribution}

$V_{\mathrm{geom}}(q)$ is a \textbf{single-well} potential with a unique
minimum at $q = 0$ (Steiner). It does not, by itself, produce a double-well
or a barrier in the tunneling sense. What $V_{\mathrm{geom}}$ does provide
is the \textbf{geometric energy cost} for displacing the junction node from
its equilibrium. Any transition from the anchor minimum ($q = 0$) to a
displaced configuration ($q > 0$) must pay this geometric cost.

\subsection{Full Double-Well Requires Additional Terms}

The full effective potential along the chosen path is:
\begin{equation}
    V(q) = V_{\mathrm{geom}}(q) + V_{\mathrm{node}}(q) + V_{\mathrm{bulk}}(q) + \ldots
    \label{eq:V_decomposition}
\end{equation}
where $V_{\mathrm{node}}$ includes junction node energy contributions and
$V_{\mathrm{bulk}}$ includes bulk gravitational and warp-factor effects.
These additional terms are \emph{not computed} in this appendix---they
require completing Put~C steps C2--C4.

For the full $V(q)$ to have double-well structure (a secondary minimum at
$q_n > 0$), the additional terms must create a sufficiently deep attractive
region at $q_n$ to overcome the geometric cost $\Delta V_{\mathrm{geom}}(q_n) > 0$.

\subsection{Conditional Barrier Statement}

\begin{resultbox}[title=Conditional Barrier Existence]
    \textbf{If} the full $V(q)$ has two local minima at $q = 0$ and
    $q = q_n > 0$ (anchor and metastable), \textbf{then} by continuity
    there exists a local maximum (barrier) at some $q_B \in (0, q_n)$.

    The geometric contribution $V_{\mathrm{geom}}(q_B) > V_{\mathrm{geom}}(0)$
    provides a \textbf{lower bound} on the barrier height: the barrier cannot
    be lower than the geometric arm-stretching cost at $q_B$.

    \textbf{Status:} Barrier existence is conditional on the double-well
    structure of $V(q)$, which remains [P] pending Put~C completion. \tagDc
\end{resultbox}

% ============================================================================
\section{Secondary Minimum Status}
\label{sec:vq:minimum}
% ============================================================================

The secondary minimum at $q = q_n$ is \textbf{postulated}, not derived:

\begin{itemize}
    \item No Hessian computation exists at $q_n$ in any donor file.
    \item The stability argument ($a(q_n) > 0$ in the Landau expansion)
          rests on observational absence of doublet partners [BL]+[M],
          not on explicit computation.
    \item $V_{\mathrm{geom}}$ alone does not produce a secondary minimum.
    \item The additional terms ($V_{\mathrm{node}}, V_{\mathrm{bulk}}$)
          that could create a minimum at $q_n$ are not computed.
\end{itemize}

\textbf{Current status:} Secondary minimum = [P], with [BL]+[M]
observational support. Promotion to [Dc] or [P*] would require at least
computing $V''(q_n) > 0$ along the chosen path from the full effective
potential, which in turn requires the non-geometric terms in
Eq.~\eqref{eq:V_decomposition}.

% ============================================================================
\section{Barrier Height Status}
\label{sec:vq:height}
% ============================================================================

Three distinct questions must be separated:

\subsection{Question 1: Barrier Existence}

Conditional on double-well structure. See \S\ref{sec:vq:barrier} above.
\textbf{Status:} [Dc] (conditional).

\subsection{Question 2: Barrier Height in Geometric Units}

The $\Zthree$ symmetry argument
(Chapter~\ref{ch:metastable}~\S\ref{sec:metastable:barrier}) gives:
\begin{align}
    E_B &= 3 \times \Earm & &\text{(barrier: all arms equally stressed)} \tagDer \\
    E_n &= 1 \times \Earm & &\text{(metastable: one effective mode)} \tagP \\
    \VB &= E_B - E_n = 2\Earm & &\tagP
    \label{eq:VB_2Earm_app}
\end{align}

The factor of 3 at the barrier is [Der] from $\Zthree$ symmetry and energy
additivity (Chapter~\ref{ch:metastable}, Step~1). The assignment $E_n = 1 \times \Earm$
(one effective deformation mode for the metastable state) is a
\textbf{model-internal postulate} [P]: it depends on the interpretation
of the metastable junction as a single-mode excitation.

Therefore $\VB = 2\Earm$ is [P]---it follows from the $\Zthree$ factor
[Der] and the single-mode assignment [P].

\subsection{Question 3: $\VB = 2\Delta m_{np}$}

The identification $\Earm \equiv \Delta m_{np}$ requires mapping the
geometric per-arm deformation energy to the observer-frame mass difference.
This mapping depends on:
\begin{itemize}
    \item the explicit form of $V(q)$ from the 5D action (Put~C completion),
    \item the projection from 5D energy to observer-frame mass.
\end{itemize}

Neither is available. \textbf{Status:} $\VB = 2\Delta m_{np}$ remains
[OPEN conditional]. The numerical agreement
($\VB^{\mathrm{cal}} \approx 2.6\MeV$ vs.\ $2\Delta m_{np} \approx 2.59\MeV$,
within 0.5\%) is a consistency check [Check], not derivation evidence.

% ============================================================================
\section{Effective Mass Status}
\label{sec:vq:mass}
% ============================================================================

The effective mass $M(q)$ in $S_{\mathrm{eff}}[q] = \int dt\,[\frac{1}{2}
M(q)\dot{q}^2 - V(q)]$ arises from the kinetic energy of junction node
displacement (Chapter~\ref{ch:instanton}).

\textbf{Scaling estimate:} Dimensional analysis within the Nambu--Goto
framework gives
\begin{equation}
    M \sim \frac{\tau R}{c^2}
    \label{eq:M_scaling}
\end{equation}
where $\tau R$ has dimensions of energy and $c$ is the speed of light. This
estimate reflects the inertia of displacing the node against arm tension.
The numerical coefficient depends on the junction geometry, path choice,
and contributions from bulk modes, and is \textbf{undetermined}.

\textbf{$q$-dependence:} $M(q) = M$ (constant) is assumed, not verified.
The q-dependence of $M$ would require computing the full kinetic term from
the 5D action, which is part of Put~C step C4.

\textbf{Status:} $M(q)$ scaling = [P] (dimensional analysis). Numerical
coefficient = [P]. $q$-independence = [P]. No donor content exists for any
of these.

% ============================================================================
\section{Oscillation Frequency at Anchor Minimum}
\label{sec:vq:omega}
% ============================================================================

From the geometric curvature (Theorem~\ref{thm:steiner_min}) and the
scaling mass estimate:
\begin{equation}
    \omega_{\mathrm{geom}}^2 = \frac{V_{\mathrm{geom}}''(0)}{M}
    = \frac{3\tau/(2R)}{M}
    \label{eq:omega_geom}
\end{equation}

This is the oscillation frequency at the anchor minimum due to the geometric
restoring force alone. The full oscillation frequency $\omegaz$ at the
metastable minimum (used in the instanton calculation,
Chapter~\ref{ch:instanton}) depends on $V''(q_n)$ of the full potential
and may differ significantly.

\textbf{Status:} $\omega_{\mathrm{geom}}$ = [Dc] within the chosen-path
model. The physical $\omegaz$ = [P].

% ============================================================================
\section{What Put~C Would Additionally Provide}
\label{sec:vq:putc}
% ============================================================================

Completing Put~C steps C2--C4 would upgrade the following:

\begin{center}
\begin{tabular}{lll}
\toprule
\textbf{Quantity} & \textbf{Current (post-R3)} & \textbf{After Put~C} \\
\midrule
$V_{\mathrm{geom}}(q)$ & [Dc] (this appendix) & subsumed by full $V(q)$ \\
$V_{\mathrm{node}}(q) + V_{\mathrm{bulk}}(q)$ & [OPEN] & [Der] \\
Full $V(q)$ & [P] (double-well postulated) & [Der] \\
Secondary minimum at $q_n$ & [P] & [Der] (if exists) \\
$\VB$ in geometric units & [P] ($2\Earm$) & [Der] \\
$M(q)$ & [P] (scaling) & [Der] \\
$\Earm$ identification & [OPEN] & [Der] \\
\bottomrule
\end{tabular}
\end{center}

\noindent Put~C is the next true 5D step. This appendix does \emph{not}
complete any part of Put~C; it bypasses C2--C4 by computing the Nambu--Goto
arm-length contribution directly.

% ============================================================================
\section{Summary}
\label{sec:vq:summary}
% ============================================================================

\begin{resultbox}[title=R3 Partial Closure Summary]
    \textbf{What R3 established:}
    \begin{itemize}
        \item $V_{\mathrm{geom}}(q) = \tau \times L_{\mathrm{tot}}(q)$
              computed explicitly along two chosen paths \hfill [Dc]
        \item Steiner minimum at $q = 0$: $V_{\mathrm{geom}}'(0) = 0$,
              $V_{\mathrm{geom}}''(0) = 3\tau/(2R) > 0$ \hfill [Dc]
        \item $V_{\mathrm{geom}}$ is a single well
              (no secondary minimum from geometry alone) \hfill [Dc]
        \item If double-well exists, $V_{\mathrm{geom}}$ provides a
              geometric lower bound on barrier cost \hfill [Dc]
        \item Factor of 3 at $\Zthree$-symmetric barrier:
              $E_B = 3\Earm$ \hfill [Der]
    \end{itemize}

    \textbf{What remains [P] or [OPEN]:}
    \begin{itemize}
        \item Double-well structure of full $V(q)$ \hfill [P]
        \item Secondary minimum existence and stability \hfill [P]
        \item $\VB = 2\Earm$ (depends on $E_n = \Earm$ assignment) \hfill [P]
        \item $\Earm \equiv \Delta m_{np}$ (unit identification) \hfill [OPEN]
        \item $M(q)$ (scaling only, no coefficient or $q$-dependence) \hfill [P]
        \item Put~C steps C2--C4 \hfill [OPEN]
    \end{itemize}

    \textbf{What this appendix does NOT claim:}
    \begin{itemize}
        \item Unique derivation of $V(q)$ from the 5D action
        \item Completion of any Put~C step
        \item Closure of barrier height in observer-frame units
        \item Derivation of $M(q)$ beyond dimensional scaling
    \end{itemize}
\end{resultbox}
